% مشقوں کا ذخیرہ — باب 4
% باب کا تعین اس فائل کے نام سے ہوتا ہے۔

\begin{exo}[کسی راستے کے ساتھ تفرقی صورت کا تکمل]{integrale-forme-differentielle}
\keywords{تفرقی صورت، کامل تفرق، شوارز کا معیار، خطی تکمل}

دو تفرقی صورتوں پر غور کریں:
\[
\omega_1 = y\,\mathrm{d}x + x\,\mathrm{d}y,
\qquad
\omega_2 = y\,\mathrm{d}x,
\]
اور مبدا $O(0,0)$ کو نقطہ $M(1,1)$ سے ملانے والے تین راستے: راستہ
$\gamma_1$ افقی قطعہ $O\to(1,0)$ اور اس کے بعد عمودی قطعہ $(1,0)\to M$
پر مشتمل ہے؛ راستہ $\gamma_2$ عمودی قطعہ $O\to(0,1)$ اور پھر افقی قطعہ
$(0,1)\to M$ پر مشتمل ہے؛ اور راستہ $\gamma_3$ وہ سیدھا قطعہ ہے جو $O$ کو
براہ راست $M$ سے ملاتا ہے۔

\begin{enumerate}
\item ایسا تفاعل $f$ پیش کر کے دکھائیں کہ $\omega_1$ کامل ہے اور $\omega_1=\mathrm{d}f$۔
\item ہر قطعے کی پیرامیٹرائزیشن سے $\int_{\gamma_1}\omega_1$ اور $\int_{\gamma_2}\omega_1$ معلوم کریں۔ حساب کے بغیر بھی نتیجہ اخذ کریں۔
\item $\omega_2$ پر شوارز کا معیار لگائیں۔ کیا نتیجہ نکلتا ہے؟
\item $\int_{\gamma_1}\omega_2$، $\int_{\gamma_2}\omega_2$ اور $\int_{\gamma_3}\omega_2$ معلوم کر کے تبصرہ کریں۔
\end{enumerate}

\begin{indication}
افقی قطعے پر $\mathrm{d}y=0$ اور عمودی قطعے پر $\mathrm{d}x=0$ ہوتا ہے؛
یوں ہر تکمل ایک سادہ تکمل بن جاتا ہے۔
\end{indication}

\begin{solution}
\textbf{سوال 1۔} $f(x,y)=xy$ موزوں ہے، کیونکہ $\partial f/\partial x=y$ اور $\partial f/\partial y=x$۔

\textbf{سوال 2۔} $t\mapsto\bigl(x(t),y(t)\bigr)$ سے پیرامیٹرائز منحنی کے لیے
\[
\int_\gamma\omega_1
= \int \left[y(t)x'(t)+x(t)y'(t)\right]\,\mathrm{d}t.
\]

راستہ $\gamma_1$ دو قطعوں کا مجموعہ ہے۔ پہلے پر $t\in[0,1]$ کے لیے
$x(t)=t$ اور $y(t)=0$؛ لہٰذا $\mathrm{d}x=\mathrm{d}t$ اور
$\mathrm{d}y=0$۔ دوسرے پر پھر $t\in[0,1]$ کے لیے $x(t)=1$ اور $y(t)=t$؛
لہٰذا $\mathrm{d}x=0$ اور $\mathrm{d}y=\mathrm{d}t$۔ پس
\[
\begin{aligned}
I_1
= \int_{\gamma_1}\omega_1
&= \int_0^1\left(0\times 1+t\times 0\right)\,\mathrm{d}t
 + \int_0^1\left(t\times 0+1\times 1\right)\,\mathrm{d}t \\
&= 0+\int_0^1\mathrm{d}t = 1.
\end{aligned}
\]

$\gamma_2$ میں پہلے قطعے کی پیرامیٹرائزیشن $x(t)=0$، $y(t)=t$ اور دوسرے کی
$x(t)=t$، $y(t)=1$ ہے، دونوں میں $t\in[0,1]$۔ چنانچہ
\[
\begin{aligned}
I_2
= \int_{\gamma_2}\omega_1
&= \int_0^1\left(t\times 0+0\times 1\right)\,\mathrm{d}t
 + \int_0^1\left(1\times 1+t\times 0\right)\,\mathrm{d}t \\
&= 0+\int_0^1\mathrm{d}t = 1.
\end{aligned}
\]
دونوں قیمتیں برابر ہیں۔ حساب کے بغیر بھی یہ معلوم تھا: چونکہ $f(x,y)=xy$
کے لیے $\omega_1=\mathrm{d}f$، اس کا تکمل صرف آخری سروں پر منحصر ہے،
\[
\int_\gamma\omega_1=f(M)-f(O)=f(1,1)-f(0,0)=1.
\]

\textbf{سوال 3۔} لکھیں $\omega_2=A(x,y)\,\mathrm{d}x+B(x,y)\,\mathrm{d}y$۔
یہاں $A(x,y)=y$ اور $B(x,y)=0$۔ اگر صورت کامل ہوتی تو شوارز کا معیار دیتا
\[
\frac{\partial A}{\partial y}=\frac{\partial B}{\partial x}.
\]
لیکن $\partial A/\partial y=1$ جبکہ $\partial B/\partial x=0$۔ مخلوط مشتق
مختلف ہیں، اس لیے $\omega_2$ کامل نہیں۔ اس کا تکمل $O$ اور $M$ کے درمیان
اختیار کیے گئے راستے پر منحصر ہو سکتا ہے، جیسا اگلا حساب دکھاتا ہے۔

\textbf{سوال 4۔} چونکہ $\omega_2=y\,\mathrm{d}x$، اس لیے صرف ایسی $x$ کی تبدیلی
تکمل میں حصہ ڈال سکتی ہے جو $y$ کی غیر صفر قیمت پر ہو۔ سابقہ پیرامیٹرائزیشن
سے $\gamma_1$ پر
\[
\begin{aligned}
J_1
=\int_{\gamma_1}\omega_2
&=\int_0^1 0\times 1\,\mathrm{d}t
  +\int_0^1 t\times 0\,\mathrm{d}t \\
&=0.
\end{aligned}
\]
$\gamma_2$ کے پہلے قطعے پر $x(t)=0$، لہٰذا $\mathrm{d}x=0$؛ دوسرے پر
$x(t)=t$، $y(t)=1$ اور $\mathrm{d}x=\mathrm{d}t$۔ پس
\[
\begin{aligned}
J_2
=\int_{\gamma_2}\omega_2
&=\int_0^1 t\times 0\,\mathrm{d}t
  +\int_0^1 1\times 1\,\mathrm{d}t \\
&=1.
\end{aligned}
\]
آخر میں قطری قطعہ $\gamma_3$ واضح طور پر یوں پیرامیٹرائز ہے:
\[
x(t)=t,\qquad y(t)=t,\qquad t\in[0,1],
\]
اس لیے $\mathrm{d}x=\mathrm{d}t$ اور $\mathrm{d}y=\mathrm{d}t$۔ نتیجتاً
\[
J_3
=\int_{\gamma_3}\omega_2
=\int_0^1 y(t)x'(t)\,\mathrm{d}t
=\int_0^1 t\,\mathrm{d}t
=\frac12.
\]
تینوں راستوں کے آخری سر یکساں ہیں مگر قیمتیں مختلف ہیں: $J_1=0$، $J_2=1$
اور $J_3=1/2$۔ یہی ناکامل تفرقی صورت کی مخصوص راستے پر منحصر خاصیت ہے۔
\end{solution}
\end{exo}

\begin{exo}[اندرونی توانائی میں یکساں تبدیلی، منتقلی کے تین طریقے]{meme-delta-u-trois-transferts}
\keywords{پہلا اصول، اندرونی توانائی، حرارت، میکانی کام، برقی کام، حرارت پیمائی، حالت کا تفاعل}

ایک مائع، اس کا حرارت پیما اور اس میں ڈوبا ہوا چھوٹا مزاحم مل کر ایک بند نظام
بناتے ہیں جو کلان پیمانے پر ساکن ہے۔ ظرف صلب ہے اور نظام کی کل حرارتی گنجائش
کو مستقل فرض کیا گیا ہے:
\[
C=2{,}00\ \mathrm{kJ\,K^{-1}}.
\]
نظام کو $T_A=20{,}0\,^\circ\mathrm{C}$ پر توازنی حالت $A$ سے
$T_B=25{,}0\,^\circ\mathrm{C}$ پر توازنی حالت $B$ تک لے جانا ہے۔ مانتے ہیں کہ
\[
U(B)-U(A)=C(T_B-T_A).
\]
کلان حرکی اور وضعی توانائیوں کی تبدیلیاں صفر ہیں۔ باہر ہونے والے تمام نقصانات نظرانداز ہیں۔

ایک ہی ابتدائی حالت $A$ سے مندرجہ ذیل تین طریقۂ کار الگ الگ کیے جاتے ہیں:
\begin{itemize}
\item[طریقۂ کار 1] نظام کو کسی کام کی فراہمی کے بغیر ایک گرم منبع سے حرارتی اتصال میں رکھا جاتا ہے۔
\item[طریقۂ کار 2] دیواریں حرارت سے موصل ہیں۔ ایک بلیڈ، کمیت $m$ کے $h=5{,}00$ m بلندی سے گرنے کے سبب مائع کو ہلاتا ہے۔ آخر میں بلیڈ اور مائع ساکن ہو جاتے ہیں اور کمیت سے ضائع ہونے والی پوری میکانی توانائی نظام کو منتقل ہوتی ہے۔ $g=9{,}81\ \mathrm{m\,s^{-2}}$ لیں۔
\item[طریقۂ کار 3] نظام حرارت $Q_3=4{,}00$ kJ وصول کرتا ہے اور باقی توانائی مزاحم کو برقی طور پر دی جاتی ہے۔ وصول شدہ برقی طاقت مستقل اور $\mathcal P=300$ W ہے۔
\end{itemize}

\begin{enumerate}
\item تینوں طریقوں کے لیے مشترک اندرونی توانائی کی تبدیلی $\Delta U=U(B)-U(A)$ معلوم کریں۔
\item پہلے دو طریقوں میں سے ہر ایک کے لیے نظام کی وصول کردہ حرارت $Q$ اور کام $W$ معلوم کریں۔ دوسرے طریقے میں مطلوب کمیت $m$ نکالیں۔
\item تیسرے طریقے کے لیے وصول شدہ برقی کام اور مزاحم کو توانائی دینے کا دورانیہ معلوم کریں۔
\item تینوں تبدیلیوں کی ابتدائی اور آخری حالتیں ایک جیسی ہیں۔ پھر بھی $Q$ اور $W$ کی قدریں مختلف کیوں ہو سکتی ہیں؟ کیا آخری حالت $B$ میں نظام کے اندر ”حرارت“ یا ”کام“ موجود ہے؟
\end{enumerate}

\begin{indication}
بینکار کے علامتی اصول کے ساتھ $\Delta U=Q+W$ استعمال کریں۔ بلیڈ کا وصول کردہ
کام، کمیت کی وضعی توانائی میں $mgh$ کمی کے برابر ہے۔ تاروں کے ذریعے سرحد پار
کرنے والی برقی توانائی کو برقی کام شمار کیا جاتا ہے۔
\end{indication}

\begin{solution}
\textbf{سوال 1۔} درجہ حرارت کا فرق $T_B-T_A=5{,}0$ K ہے۔ لہٰذا
\[
\Delta U=C(T_B-T_A)
=2{,}00\times5{,}0
=10{,}0\ \mathrm{kJ}.
\]
یہ قیمت صرف توازنی حالتوں $A$ اور $B$ پر منحصر ہے۔

\textbf{سوال 2۔} پہلے طریقے میں کوئی کام وصول نہیں ہوتا: $W_1=0$۔ پہلا اصول دیتا ہے
\[
Q_1=\Delta U=10{,}0\ \mathrm{kJ}.
\]
حرارت مثبت ہے کیونکہ توانائی نظام میں داخل ہوتی ہے۔

دوسرے طریقے میں دیواریں حرارت سے موصل ہیں، اس لیے $Q_2=0$۔ اندرونی توانائی
کی پوری تبدیلی بلیڈ کے میکانی کام سے آتی ہے:
\[
W_2=\Delta U=10{,}0\ \mathrm{kJ}.
\]
چونکہ $W_2=mgh$، اس لیے
\[
m=\frac{W_2}{gh}
=\frac{10{,}0\times10^3}{9{,}81\times5{,}00}
\simeq 204\ \mathrm{kg}.
\]
یہ بڑی مقدار یاد دلاتی ہے کہ درجہ حرارت میں معمولی اضافہ بھی خاصی میکانی
توانائی کے برابر ہوتا ہے۔

\textbf{سوال 3۔} پہلا اصول لازم کرتا ہے کہ
\[
W_3=\Delta U-Q_3=10{,}0-4{,}00=6{,}00\ \mathrm{kJ}.
\]
یہ کام مثبت ہے: برقی رسد نظام کو توانائی دیتی ہے۔ $W_3=\mathcal P\,\Delta t$ سے
\[
\Delta t=\frac{W_3}{\mathcal P}
=\frac{6{,}00\times10^3}{300}
=20{,}0\ \mathrm{s}.
\]

\textbf{سوال 4۔} تینوں راستے بالترتیب دیتے ہیں
\[
(Q_1,W_1)=(10{,}0,0)\ \mathrm{kJ},
\qquad
(Q_2,W_2)=(0,10{,}0)\ \mathrm{kJ},
\]
اور
\[
(Q_3,W_3)=(4{,}00,6{,}00)\ \mathrm{kJ}.
\]
ہر صورت میں مجموعہ $Q+W$، $10{,}0$ kJ ہے اور یکساں تبدیلی $\Delta U$ پیدا
کرتا ہے۔ اندرونی توانائی حالت کا تفاعل ہے، جبکہ حرارت اور کام کسی تبدیلی کے
دوران منتقلیوں کو بیان کرتے ہیں۔ آخری حالت $B$ میں نظام اندرونی توانائی رکھتا
ہے، مگر اس میں نہ ”حرارت“ ہے نہ ”کام“۔
\end{solution}
\end{exo}

\begin{exo}[مستقل بیرونی دباؤ کے تحت انضغاط]{compression-pression-exterieure-constante}
\seoready{true}
\keywords{دباؤ کی قوتوں کا کام، مستقل بیرونی دباؤ، انضغاط، پھیلاؤ، خلا میں پھیلاؤ}

ایک گیس کو مستقل بیرونی دباؤ $P_0$ کے تحت حجم $V_A$ سے حجم $V_B<V_A$ تک دبایا جاتا ہے۔

\begin{enumerate}
\item گیس کے وصول کردہ کام کا حساب کریں۔
\item گیس اسی بیرونی دباؤ $P_0$ کے خلاف $V_B$ سے $V_A$ تک لوٹتی ہے۔ وصول کردہ کام معلوم کر کے انضغاط کے کام سے موازنہ کریں۔
\item جب گیس خلا میں پھیلے تو $V_B$ سے $V_A$ تک پھیلاؤ دوبارہ حل کریں۔ حاصل شدہ تینوں علامتوں کی تشریح کریں۔
\end{enumerate}

\begin{indication}
$W=-\int P_{\mathrm{ext}}\,\mathrm{d}V$ استعمال کریں۔ تبدیلی نیم ساکن نہ بھی ہو
تو \emph{بیرونی} دباؤ ہی کا تکمل لینا ہے۔
\end{indication}

\begin{solution}
\textbf{سوال 1۔} بیرونی دباؤ مستقل ہے، اس لیے
\[
W_{A\to B}=-P_0(V_B-V_A)=P_0(V_A-V_B)>0.
\]
انضغاط کے دوران گیس کام وصول کرتی ہے۔

\textbf{سوال 2۔} واپسی کے لیے
\[
W_{B\to A}=-P_0(V_A-V_B)<0.
\]
دونوں کام ایک دوسرے کے مخالف ہیں کیونکہ دونوں تبدیلیاں اسی مستقل بیرونی دباؤ کے خلاف کی گئی ہیں۔

\textbf{سوال 3۔} خلا میں $P_{\mathrm{ext}}=0$، لہٰذا
\[
W_{B\to A}^{\mathrm{vide}}=0.
\]
اس لیے پھیلاؤ کا مطلب خودبخود منفی کام نہیں: کسی بیرونی قوت کو حقیقت میں پیچھے دھکیلنا ضروری ہے۔
\end{solution}
\end{exo}
